% 来源：OEBPS/Text/01_10.xhtml（OPF spine 顺序；仅第一册）
\renewcommand{\bookchapter}{4}
\cleardoublepage
\BookTOCChapter[第四章\q 芒格主义：查理的即席谈话]{第四章\\ 芒格主义：查理的即席谈话}

\vspace{-0.8em}

{\centering \textbf{——在伯克希尔·哈撒韦公司和西科金融公司年会上的谈话摘录}\par}

\vspace{1em}

沃伦·巴菲特和查理·芒格毫无疑问是有史以来最伟大的投资二人组，所以如果有投资者不去尽可能多地了解这两个人，尽可能多地学习他们如何取得成功，那么这个投资者就是——用查理很喜爱的一个词来说——“疯子”。
\par

但是向沃伦和查理学习的真正乐趣并不在于人们能够了解大量关于如何在长时间内让金钱以很高的复利增长的知识——尽管这当然是一种很不错的额外收获——而是在于通过听取他们的教导，人们将能够更加深刻地理解人类的本性、世界的现状、如何理性思考，以及最重要的，如何更好地过上一种正直、幸福、善良的生活（提示：这三种要素是相辅相成的）。向沃伦学习并不难——写他的书有很多本；他经常发表演讲，撰写文章，出席公众场合；几十年来，他每年都会发表一篇很长的信件（你可以在伯克希尔·哈撒韦网站上免费阅读最近的30封信，网址是www.berkshirehathaway.com。如果你还没去看过，那你还在等什么呢？！）。但查理就神秘多了：关于他的书只有两本。和沃伦相比，他发表的演讲和撰写的文章要少得多。
\par

由于这个原因，许多人并不知道查理本人是个天才，而且他对沃伦的投资哲学产生了深刻的影响，这一点沃伦是公开承认的。他们之间的互动会让观者忍俊不禁。在伯克希尔年会的主席台上，沃伦一般会率先回答提问，但然后往往会转过身说：“查理，你怎么看？”查理的典型反应是一动不动、面无表情地（人们很容易用一个假人来代替他，在大多数情况下保管没有人会发现）说：“我没有什么要补充的。”这九个字已经变成了查理的标志，通常让沃伦和听众会心一笑。（实际上，1998年9月，在为收购通用再保险公司而召开的特别大会上，沃伦真的用一个硬纸板做成的查理人像和一段查理说“我没有什么要补充的”的录音代替查理出现在主席台上。我的一个朋友参加了该大会，他说：“沃伦大概播了六次这段录音，每次都会露出顽皮的神色。”）这几个字也反映了查理古怪的脾气——至少他的公众形象是这样的——以及这两个心灵相通的杰出生意人“合二为一”的关系。
\par

但每当查理确实有话要补充，那么他说的话往往很尖锐、很深刻，而且他也不会拐弯抹角。用他自己的话来说，查理是一个“古怪的、老派的”人，而且从来不为政治正确而担忧——他想到什么就说什么。
\par

因为在伯克希尔年会上，大多数讲话都是由沃伦来发表，所以我总是喜欢参加几天后在帕萨迪纳举行的西科年会，去领略查理深邃的思想（他是西科金融公司的主席，西科是一家结构和伯克希尔相同的控股公司，后者拥有西科80.1\%的股份）。人们不用是股东也可以参加。查理——就像他和沃伦在伯克希尔年会上那样——欢迎所有来到格雷厄姆、多德、巴菲特和芒格教堂的朝圣者。我从来不曾后悔从奥马哈回来没几天就长途奔波，横跨大半个美国去听查理演讲，因为他的思维依然处在巅峰状态。在2003年的会议上，查理做出一个又一个充满睿见的点评，我有个朋友侧过身来，低声对我说：“真难相信他已经80岁了！”确实如此！
\par

作为导师和思想家，查理同沃伦旗鼓相当，我从他那里学到了很多知识。而在这里，我把自己过去五年中在伯克希尔和西科年会上所作的记录抽取出来，包括了查理在这些会议上说过的最有洞察力、最发人深思、最抒发胸臆、最好玩有趣的话。由于他的大多数发言都是为了回应股东的提问，所以他广泛地谈论了伯克希尔·哈撒韦的情况，还有许多和商业、投资相关的话题。
\par

为了让查理的教诲更加容易被接受，我依照话题对它们进行了组织，也稍微做了些编辑工作，但目的是让他亲自说出他的话。请尽情享受吧！
\par

\begin{WisdomBox}
\WisdomSpeaker{查理·芒格}
\par

我和一个眼光极准的传奇人物合作了许多年，我总是很惊讶地发现，原来人们只要掌握和不停地应用所有明显而且容易学习的原理，就能够精通许许多多的领域。
\par
\end{WisdomBox}

\subsection*{芒格之所以成功的两个关键因素}

1. 如果乌龟能够吸取它那些最棒前辈的已经被实践所证明的洞见，有时候它也能跑赢那些追求独创性的兔子，或者跑赢宁愿跻身那些忽略前人最优秀的工作的蠢货之列的兔子。乌龟若能找到某些特别有效的方法来应用前人最伟大的工作，或者只要能避免犯下常见的错误，这种情况就会发生。我们赚钱，靠的是记住浅显的，而不是掌握深奥的。我们从来不去试图成为非常聪明的人，而是持续地试图别变成蠢货，久而久之，我们这种人便能获得非常大的优势。
\par

\begin{WisdomBox}
\WisdomSpeaker{查理·芒格}
\par

我们从来不去试图成为非常聪明的人，而是持续地试图别变成蠢货。
\par
\end{WisdomBox}

2. 我们并不自称是道德高尚的人，但至少有很多即便合法的事情，也是我们不屑去做的。我们不会去做那些事情。目前美国有种文化认为，所有不会把你送进监狱的事情都是可以做的事情。
\par

我们认为，在你应该做的事情和就算你做了也不会受到法律制裁的事情之间应该有一条巨大的鸿沟。我想你应该远离那条线。我觉得我们不应该由此而得到太多的赞誉。这样做事的原则帮我们赚到更多的钱。我相信就算这种经营方式没有给我们赚这么多钱，我们也不会做坏事。但更多的时候，我们由于做正确的事情而赚到更多的钱。
\par

记住路易斯·文森狄的规则：说真话，你将无须记住你的谎言。就是这么简单的一个概念。
\par

\begin{WisdomBox}
\WisdomSpeaker{查理·芒格}
\par

在你应该做的事情和就算你做了也不会受到法律制裁的事情之间应该有一条巨大的鸿沟。
\par
\end{WisdomBox}

\subsection*{如何经营好自己的生活}

\subsubsection*{如何获得幸福和成功}

如果你在生活中唯一的成功就是通过买股票发财，那么这是一种失败的生活。生活不仅仅是精明地积累财富。
\par

生活和生意上的大多数成功来自你知道应该避免哪些事情：过早死亡、糟糕的婚姻等等。
\par

避免染上艾滋病、在路口和火车抢道，以及吸毒。培养良好的心理习惯。
\par

避免邪恶之人，尤其是那些性感诱人的异性。
\par

如果因为你的特立独行而在周围人中不受欢迎……那你就随他们去吧。
\par

\begin{WisdomBox}
\WisdomSpeaker{查理·芒格}
\par

生活不仅仅是精明地积累财富。
\par
\end{WisdomBox}

\subsubsection*{满足于你已经拥有的}

这里是一句你的投资顾问可能会反对的真话：如果你已经相当富裕，而别人的财富增长速度比你更快，比如说，在投资高风险的股票上赚了钱，那又怎样呢？！总是会有人的财富增长速度比你快。这并不可悲。
\par

看看斯坦利·德鲁肯米勒（他管理乔治·索罗斯的基金，由于投机高科技股票和生物科技股票而蒙受惨重的损失）吧：他总是必须成为最好的，无法容忍其他人在这些投资领域击败他。
\par

索罗斯无法忍受其他人从科技产业赚钱而自己没有赚到，他亏得一塌糊涂。我们根本就不在意（别人在科技产业赚了钱）。
\par

\subsubsection*{提防妒忌}

关注别人赚钱（比你）更快的想法是一种致命的罪行。妒忌真的是一种愚蠢的罪行，因为它是仅有的一种你不可能得到任何乐趣的罪行。它只会让你痛苦不堪，不会给你带来任何乐趣。你为什么要妒忌呢？
\par

\subsubsection*{如何致富}

（有个年轻的股东问查理要怎样才能追随他的足迹，查理的回答折服了在场的所有人。）
\par

许多年轻的生意人都问我们这个问题。这是一个聪明的问题：你看见某个有钱的老头，然后问他：“我要怎样才能变得像你一样呢？”
\par

每天起床的时候，争取变得比你从前更聪明一点。认真地、出色地完成你的任务。慢慢地，你会有所进步，但这种进步不一定很快。但你这样能够为快速进步打好基础……每天慢慢向前挪一点。到最后——如果你足够长寿的话——大多数人得到了他们应得的东西。
\par

\begin{WisdomBox}
\WisdomSpeaker{查理·芒格}
\par

每天起床的时候，争取变得比你从前更聪明一点。
\par
\end{WisdomBox}

\subsubsection*{如何找到好配偶}

找到好配偶的最佳方法是什么呢？唯一的最佳方法是你首先必须值得拥有好配偶，因为从定义上来说，好配偶可绝不是傻瓜。
\par

\subsubsection*{阅读的重要性}

我这辈子遇到的聪明人（来自各行各业的聪明人）没有不每天阅读的——没有，一个都没有。沃伦读书之多，我读书之多，可能会让你感到吃惊。我的孩子们都笑话我。他们觉得我是一本长了两条腿的书。
\par

我本人是个传记书迷。我觉得你要是想让人们认识有用的伟大概念，最好是将这些概念和提出它们的伟人的生活和个性联系起来。我想你要是能够和亚当·斯密交朋友，那你的经济学肯定可以学得更好。和“已逝的伟人”交朋友，这听起来很好玩，但如果你确实在生活中与已逝的伟人成为朋友，那么我认为你会过上更好的生活，得到更好的教育。这种方法比简单地给出一些基本概念好得多。
\par

\begin{WisdomBox}
\WisdomSpeaker{查理·芒格}
\par

如果你确实在生活中与已逝的伟人成为朋友，那么我认为你会过上更好的生活，得到更好的教育。
\par
\end{WisdomBox}

\subsubsection*{减少物质需求}

大多数人将会发现回报率下降了（由于通货膨胀）。如果你担心通货膨胀，最好的预防手段之一就是在你的生活中别拥有大量愚蠢的需求——你不需要很多物质的商品。
\par

\begin{WisdomBox}
\WisdomSpeaker{查理·芒格}
\par

如果你担心通货膨胀，最好的预防手段之一就是在你的生活中别拥有大量愚蠢的需求——你不需要很多物质的商品。
\par
\end{WisdomBox}

\subsubsection*{慈善}

（沃伦和我）觉得我们这些非常幸运的人有责任反馈社会。至于是像我这样在活着的时候捐赠大量的钱，还是像沃伦那样，先捐一点，（等去世之后）再捐很多，那是个人喜好的问题。我恐怕不会喜欢人们整天问我要钱。沃伦更无法忍受这种情况。
\par

\subsubsection*{避免欠债}

你一旦走进负债的怪圈，就很难走出来。千万别欠信用卡的钱。你无法在支付18\%的利息的情况下取得进展。
\par

\subsubsection*{公立学校的衰落}

你完全可以说这（公立学校的衰落）是我们时代的主要灾难之一。我们将地球史上最大的成功之一变成了地球史上最大的灾难之一。
\par

\subsubsection*{日本的经济衰退}

日本的经济衰退足以令每个人大惊失色，它持续了十年，尽管利率低于1\%。政府把所有的货币政策游戏都玩过了，但毫无用处。如果你向哈佛大学的经济学家描绘这种情况，他们会说这是不可能的。然而与此同时，香港却出现了资产泡沫。为什么？因为日本和中国拥有两种非常不同的文化。中国人喜欢赌博。
\par

这个经典的例子说明了为什么成功的投资者必须吸取各个学科的知识。想象一下有个经济学家在经济学研讨大会上站起来说出我的解释。那将会是政治不正确！但经济学的工具无法解释现实的情况。
\par

\subsection*{如何像芒格一样聪明}

\subsubsection*{多元思维模型的重要性}

你必须知道重要学科的重要理论，并经常使用它们——要全部都用上，而不是只用几种。大多数人都只养成一个学科——比如说经济学——的思维模型，试图用一种方法来解决所有问题。你知道谚语是怎么说的：“在手里拿着铁锤的人看来，世界就像一颗钉子。”这是一种处理问题的笨方法。
\par

遇到不同的公司，你需要检查不同的因素，应用不同的思维模型。我无法简单地说：“就是这三点。”你必须自己把它（投资技巧）总结出来，一辈子深种在你的大脑中。
\par

\begin{WisdomBox}
\WisdomSpeaker{查理·芒格}
\par

你必须知道重要学科的重要理论，并经常使用它们。
\par
\end{WisdomBox}

如果你用的是许多学生使用的那种学习方法，那你将无法掌握你需要的那100种思想——那些学生学得很好，接着把知识丢回给他们的教授，拿到一个分数，然后就把知识清空，就好像放掉浴缸里的水，以便下次装更多水那样。如果这是你用来学习那100种思维模型的方法，在生活中你将会一无所获。你必须认真学习这些模型，让它们变成你随时能够使用的本领。
\par

顺便说一下，没有规则要求你不能增加一两个新的模型，即使你已经步入晚年。实际上，我现在还在不断学习。不过大多数有用的模型是在年轻时学到的。
\par

我推荐的这些精神领域是非常快乐的，没有人愿意从那里离开。离开就像砍断自己的双手。
\par

\subsubsection*{不自欺的精神}

大量的人很荒唐地死抱着错误的观念不放。凯恩斯说：“介绍新观念倒不是很难，难的是清除那些旧观念。”
\par

不自欺的精神是你能拥有的最好的精神。它非常强大，因为它太少见了。
\par

\begin{WisdomBox}
\WisdomSpeaker{查理·芒格}
\par

不自欺的精神是你能拥有的最好的精神。它非常强大，因为它太少见了。
\par
\end{WisdomBox}

\subsubsection*{常（非常）识}

得到系统的常（非常）识——非常基础的知识——是一种威力巨大的工具。电脑带来的危险非常多。人们计算得太多，思考得太少。
\par

拥有常识不但意味着有能力辨认智慧，也意味着有能力拒绝愚蠢。如果排除了许多事情，你就不会把自己搞得一团糟。
\par

\begin{WisdomBox}
\WisdomSpeaker{查理·芒格}
\par

拥有常识不但意味着有能力辨认智慧，也意味着有能力拒绝愚蠢。
\par
\end{WisdomBox}

我们看很多书。我认识的聪明人没有不看很多书的。但光看书还不够：你必须拥有一种能够掌握思想和做合理事情的性格。大多数人无法掌握正确的思想，或者不知道该怎么应用它们。
\par

你拥有的基本知识越多，你需要吸取的新知识就越少。那个蒙上眼睛下棋的家伙（他是个象棋大师，在伯克希尔·哈撒韦开年会那个周末赶到奥马哈展示他的棋技，蒙上眼睛同时跟好几个人对弈）——他拥有关于棋盘的知识，这使他能够那么做。
\par

\begin{WisdomBox}
\WisdomSpeaker{查理·芒格}
\par

你必须拥有一种能够掌握思想和做合理事情的性格。
\par
\end{WisdomBox}

\subsubsection*{各学科间的相互割据和相互孤立是致命的}

（美国的大学）有很多毛病。要是我的话，我会砍掉四分之三的院系——除了自然科学，一个都不留。但没有人会这么做，所以我们不得不忍受那些毛病。（高等教育）错误程度让人吃惊。学科之间毫无关联，这是致命的。大学里每个系都有些看不见大环境的古怪家伙。
\par

我认为美国各所主要大学的人文社科教授的观点都经不起推敲，至少他们的公共政策理论是这样的。（不过呢，）他们可能对法语了解蛮多的。
\par

怀特海曾经指出，学科各自孤立的情况是致命的，每个教授并不了解其他学科的思维模型，将其他学科和他自己的学科融会贯通就更别提了。我想怀特海讨厌的这种研究方法有一个现代的名字，那个名字叫作神经研究法。
\par

\subsection*{芒格对巴菲特的评论}

人们很难相信他的业绩一年比一年好。这种情况不会永久地持续下去，但沃伦的确在进步。这是很罕见的：绝大多数人到了72岁就停滞不前了，但沃伦依然在进步。
\par

\subsubsection*{查理对沃伦的影响}

我认为那些作者言过其实地夸大了我对沃伦的影响。沃伦确实有过发蒙的时候，因为他曾在本杰明·格雷厄姆手下工作过，而且赚了那么多的钱。从如此成功的经验中跳出来确实很难。但如果世上未曾有过查理·芒格这个人，巴菲特的业绩依然会像现在这么漂亮。
\par

我想那些认为我是沃伦的伟大启蒙者的想法里有好些神话的成分。他不需要什么启蒙。但我们现在比五年前更能明白这个道理。
\par

\begin{WisdomBox}
\WisdomSpeaker{查理·芒格}
\par

我想那些认为我是沃伦的伟大启蒙者的想法里有好些神话的成分。
\par
\end{WisdomBox}

\subsubsection*{沃伦走后会怎样？}

关键是拥有许多优秀的企业。优秀的企业能给伯克希尔的发展带来很多动力。（然而，）我认为我们的继任者在资本配置方面将不会像沃伦这么出色。
\par

伯克希尔钱多成灾——我们拥有许多不断产生现金的伟大企业。如果股价下跌，伯克希尔可以把它买回来。没有理由认为它会很快完蛋，而且我觉得有理由认为它将会继续运转良好。如果伯克希尔没有随着时间的流逝——就算是在沃伦去世之后——而变得更大、更强，我会感到非常吃惊。
\par

等到沃伦离开的时候，伯克希尔的收购业务会受到影响，但其他部门将会运转如常。收购业务应该也还行。反正我们可以向你们保证的是，从前那种增长速度将会下降，我们可不想在这个问题上说谎。
\par

我想到那时伯克希尔的最高领导人应该没有沃伦那么聪明。但别抱怨：“天啊，给我沃伦·巴菲特40年之后，怎么能给我一个比他差的混蛋呢？”那是很愚蠢的。
\par

如果说谁有理由担忧的话，那个人应该是我，但我认识巴菲特家族几十年了，我可以告诉你：“别担心这件事。你应该感到非常幸运。”
\par

\subsubsection*{要是查理去世了呢？}

正如你能看到的，我们打算永远在这里待下去。（但就算我去世了）你拥有大量的现金，还有巴菲特坐镇总公司——还有什么好担心的呢？
\par

\subsubsection*{查理是可恶的“说不大师”}

巴菲特：你千万要学会如何校准他的答案。如果你问查理某个项目，他说“不”，那么我们会用所有的钱来投资它。如果他说“这是我听过的最愚蠢的事情”，那么我们就会适当地做一点投资。如果你能够校准他的答案，你就能获得很多智慧。
\par

沃伦曾经告诉我一个故事，他们刚认识的时候，还没有一起做生意，各自拥有一家合伙投资公司，沃伦经常给查理打电话说“我想做某件事”，并对这件事加以描绘，查理会说：“天啊，你在开玩笑吗？有这种风险和那种风险。”他们将会仔细考虑查理看到的这些风险，然而沃伦通常会说：“我想你是对的。”但有时候他也会说：“查理，你说的我都听进去了，但我还是想做这件事。”沃伦说他直到那一刻才能知道查理的真正想法，因为查理有时会说：“沃伦，如果你要做这件事，我能不能参股？”人们通常认为查理的价值在于识别风险和说“不”，但其实查理最有价值的地方，在于他能够识别那些可以参股的时机。
\par

查理明白要找到真正好的东西很难。所以，就算你90\%的时间都在说“不”，你也不会错失太多的东西。
\par

\begin{WisdomBox}
\WisdomSpeaker{巴菲特}
\par

查理明白要找到真正好的东西很难。所以，就算你90\%的时间都在说“不”，你也不会错失太多的东西。
\par
\end{WisdomBox}

\subsection*{伟大的投资者是怎样炼成的}

\subsubsection*{性格、耐心和求知欲的重要性}

（成功投资的关键因素之一就是拥有良好的）性格——大多数人总是按捺不住，或者总是担心过度。成功意味着你要非常有耐心，然而又能够在你知道该采取行动时主动出击。你吸取教训的来源越广泛，而不是仅仅从你自己的糟糕经验中吸取教训，你就能变得越好。
\par

\begin{WisdomBox}
\WisdomSpeaker{查理·芒格}
\par

成功意味着你要非常有耐心，然而又能够在你知道该采取行动时主动出击。
\par
\end{WisdomBox}

我想说说怎样才能培养那种毫不焦躁地持有股票的性情。光靠性格是不行的。你需要在很长很长的时间内拥有大量的求知欲望。
\par

你必须有浓厚的兴趣去弄明白正在发生的事情背后的原因。如果你能够长期保持这种心态，你关注现实的能力将会逐渐得到提高。如果你没有这种心态，那么即使你有很高的智商，也注定会失败。
\par

\subsubsection*{集中投资}

我们的投资风格有一个名称——集中投资，这意味着我们投资的公司有10家，而不是100家或者400家。
\par

我们的游戏是，当好项目出现时，我们必须能够认出来，因为好项目并不会经常出现。机会只眷顾有准备的人。
\par

好的投资项目很难得，所以要把钱集中投在少数几个项目上，这在我看来是很好的主意。投资界有98\%的人并不这么想。而我们一直以来都是那么做的，这给我们——也给你们——带来了许多好处。
\par

好玩的是，绝大多数大型投资机构的想法不是这样的。他们聘请了许多人，比较默克制药公司和辉瑞制药公司到底哪个价值高，分析标准普尔500指数的每只成分股，以为他们能够打败市场。你这样做是无法打败市场的。
\par

我们的投资规矩是等待好球的出现。
\par

\begin{WisdomBox}
\WisdomSpeaker{查理·芒格}
\par

我们的投资规矩是等待好球的出现。
\par
\end{WisdomBox}

如果我有机会进入某个企业工作，而那里的人用各种指数标准来衡量我，逼我把钱都投出去，聚集在周围监视我，那么我会讨厌这家企业。我会有戴上脚镣手铐的感觉。
\par

很少人借鉴我们的方法。采用集中投资策略的公司略有增加，但真正增加的是无限制使用投资顾问来建议资产如何配置、分析评估其他顾问等等做法。也许有2\%的人会走进我们的阵营，剩下98\%的人将会信奉他们听来的道理（比如说市场是完全有效的）。
\par

\subsubsection*{错误的投资课程}

Beta系数、现代组合投资理论等等——这些在我看来都没什么道理。我们想要做到的是，用低廉的价格，甚至是合理的价格，来购买那些拥有可持续竞争优势的企业。
\par

大学教授怎么可以散播这种无稽之谈（股价波动是衡量风险的尺度）呢？几十年来，我一直都在等待这种胡言乱语结束。现在这样乱讲的人少了一些，但还是有。
\par

沃伦曾经对我说：“我对学术界的总体看法可能是错的（把它看得这么糟糕），因为那些跟我打交道的人都有一套疯狂的理论。”
\par

\subsubsection*{分散投资}

认为投资应该越分散越好是一种发疯的观念。
\par

\begin{WisdomBox}
\WisdomSpeaker{查理·芒格}
\par

认为投资应该越分散越好是一种发疯的观念。
\par
\end{WisdomBox}

我们并不认为分散程度很高的投资能够产生好结果。我们认为几乎所有好投资的分散程度都是相对较低的。
\par

如果你把我们15个最好的决策剔除，我们的业绩将会非常平庸。你需要的不是大量的行动，而是极大的耐心。你必须坚持原则，等到机会来临，你就用力去抓住它们。
\par

这些年来，伯克希尔就是通过把赌注押在有把握的事情上而赚钱的。
\par

\subsubsection*{坐等投资法}

如果你因为一样东西的价值被低估而购买了它，那么当它的价格上涨到你预期的水平时，你就必须考虑把它卖掉。那很难。但是，如果你能购买几个伟大的公司，那么你就可以安坐下来啦。那是很好的事情。
\par

我们偏向于把大量的钱投在我们不用再另作决策的地方。
\par

\subsubsection*{优秀企业是什么？}

有这么两类企业：第一类每年赚12\%，你到年底可以把利润拿走。第二类每年赚12\%，但所有多余的现金必须进行再投资——它总是没有分红。这让我想起了那个卖建筑设备的家伙——他望着那些从购买新设备的客户手里吃进的二手机器，并说：“我所有的利润都在那里了，在院子里生锈。”我们讨厌那种企业。
\par

\subsubsection*{喜诗糖果：优秀企业的个案研究}

如果喜诗糖果（在我们收购它的时候）再多要10万美元（巴菲特插口说：“1万美元”），沃伦和我就会走开——我们那时就是那么蠢。
\par

伊拉·马歇尔说你们这两个家伙疯啦——有些东西值得你们花钱买的呀，比如说优质的企业和优质的人。你们低估了优质的价值啦。我们听进了他的批评，改变了我们的想法。这对每个人来说都是很好的教训：要有能力建设性地接受批评，从批评中吸取教训。如果你间接地吸取我们从喜诗糖果学到的教训，那么你将会说，伯克希尔是建立在建设性的批评之上的。
\par

\subsubsection*{错误}

伯克希尔历史上最严重的错误是坐失良机的错误。我们看到了许多好机会，却没有采取行动。这些是巨大的错误——我们为此损失了几十亿美元。到现在还在不断地犯这样的错误。我们正在改善这个缺陷，但还没有完全摆脱它。
\par

这样的错误可以分为两类：1. 什么也不做，沃伦称之为“吮吸我的大拇指”；2. 有些股票本来应该买很多，但是只买了一点。
\par

在差点犯下没有购买喜诗糖果的大错之后，我们有很多次犯了相同的错误。很显然，我们吸取教训很慢。这些机会成本并没有反映在财务报表上，但却让我们失去了几十亿美元。
\par

由于坐失良机的错误（是无形的），绝大多数人并不会注意到它们。我们会主动说出我们犯下哪些坐失良机的错误——就像我们刚才所做的那样。（刚刚在讨论的是，他们因为沃尔玛的股价上涨了一点点而没有买进，结果犯下了少赚100亿美元的错误。）
\par

\subsubsection*{当股价下跌时买进}

过去几十年来，我们经常这么做：如果某家我们喜欢的企业的股票下跌，我们会买进更多。有时候会出现一些情况，你意识到你错了，那么就退出好了。但如果你从自己的判断中发展出了正确的自信，那么就趁价格便宜多买一些吧。
\par

诱人的机会总是转瞬即逝的。
\par

\begin{WisdomBox}
\WisdomSpeaker{查理·芒格}
\par

诱人的机会总是转瞬即逝的。
\par
\end{WisdomBox}

真正好的投资机会不会经常有，也不会持续很长的时间，所以你必须做好行动的准备。要有随时行动的思想准备。
\par

\subsubsection*{小投资者的机会}

如果你今天只有很少的资本，年纪尚轻，那么你的机会比我年轻时要少一些。当年我们刚走出大萧条。当时人们认为资本主义很糟糕。
\par

20世纪20年代出现了许多弊端。当时有个笑话是这样的，一个家伙说：“我买股票是为了晚年，这种方法很灵验——不到六个月，我就觉得自己像个老头了！”
\par

你遇到的困难更大，但这并不意味着你没办法做得好——只是需要更多时间而已。但那有什么关系呢，你能活得更长久啊。
\par

如果我是你，我就会买些很小的股票，寻找一些罕见的错误定价的机会，但这类机会的确是很少。
\par

\subsubsection*{股票卖空}

卖空股票，却又看到股价因为某些因素而上涨，是很苦恼的事情。你在生活中不应该去做这么苦恼的事情。
\par

\subsubsection*{护城河和持续的竞争优势}

旧的护城河正在被填平，而新的护城河比以前更难预测，所以事情变得越来越难。
\par

\subsubsection*{学习过程}

我还没发现有谁能够很快做到（成为一名伟大的投资者）。作为投资者，沃伦·巴菲特比我第一次遇到他时好得太多了，我也是这样的。所以诀窍就在于不断学习。你应该喜欢学习的过程。
\par

\begin{WisdomBox}
\WisdomSpeaker{查理·芒格}
\par

成为伟大投资者的诀窍在于不断学习。
\par
\end{WisdomBox}

我观察沃伦几十年了。沃伦学习了很多东西，这使他得以（拓展他的能力圈，所以他能够投资像中国石油这样的企业）。
\par

如果你打算成为一名投资者，那么总会有些投资领域不是你完全了解的。但如果你努力学习，不断进步，那么你就能够保证你的投资都会产生好的结果。关键是自律、勤奋和练习。那就像打高尔夫球——你必须努力练习。
\par

如果你没有不停地学习，其他人将会超过你。
\par

\subsubsection*{能力圈及其边界}

我们接触过的事情有很多。我们有三个选项：可以；不行；太难。除非我们对某个项目特别看好，否则就会把它归为“太难”的选项。你需要做的，只是寻找一块特殊的能力领域，并把精力放在那里面就可以了。
\par

如果你确有能力，你就会非常清楚你能力圈的边界在哪里。没有边界的能力根本不能称之为能力。（如果你问自己是否过了能力圈的范围，）那这个问题本身就是答案。
\par

\begin{WisdomBox}
\WisdomSpeaker{查理·芒格}
\par

如果你确有能力，你就会非常清楚你能力圈的边界在哪里。没有边界的能力根本不能称之为能力。
\par
\end{WisdomBox}

\subsubsection*{资本成本和机会成本}

（巴菲特：查理和我不知道我们的资本成本。商学院有教，但我们持怀疑态度。我们只是用我们拥有的资本尽量去做最聪明的事情而已。我们会拿一切跟其他选择作比较。我从来没觉得那个资本成本的计算对我来说有什么意义。查理，你呢？）
\par

我也不觉得。如果你拿起最好的经济学教材，那是曼昆写的，他说人们做出决定的依据是机会成本——也就是说，真正重要的是是否有其他选择。我们决策的时候也是这样的。其他人却不这么认为——甚至还有人提出了股本成本的说法。这些人真是脑子坏掉了。
\par

很明显，关键是考虑各种成本，包括机会成本。当然，资本不是免费的。要算出你的借贷成本很容易，但理论家对股本成本的说法是很神经的。他们说，如果你获得100\%的资本回报率，那么你就不应该去投资某个产生80\%回报率的项目。这简直是疯话。
\par

\subsubsection*{预测的价值}

人们总是喜欢请人来为他们预测未来。很久很久以前，有些国王会请人来解读绵羊的内脏。假装知道未来的人总是有市场的。聆听今天的预测家的行为，就像古代国王请人看绵羊的内脏那么疯狂。这种情况总是一而再、再而三地出现。
\par

\subsubsection*{首次公开募股（IPO）}

如果能够使用我们的思维模型，那么你完全有可能找到值得购买的新股。每年发行的新股有无穷多，我坚信你总能抓住一两个机会。但普通人玩新股很容易遭受重挫，所以你要是觉得你是个天才，那就去玩吧，祝你好运。
\par

对我们来说，新股要么发行规模太小，要么太过高科技——我们不懂它们是干什么的。如果沃伦在关注它们，那我就不知道啦。
\par

\subsection*{如何洞悉金融市场}

\subsubsection*{股票、伦勃朗和泡沫}

股票有点像债券，对其价值的评估，大略以合理地预测未来产生的现金为基础。但股票也有点像伦勃朗的画作，人们购买它们，是因为它们的价格一直都在上涨。这种情况，再加上股票涨跌所产生的巨大“财富效应”，可能会造成许多祸害。
\par

\begin{WisdomBox}
\WisdomSpeaker{查理·芒格}
\par

但股票也有点像伦勃朗的画作，人们购买它们，是因为它们的价格一直都在上涨。
\par
\end{WisdomBox}

（2001年4月）这是资本主义最极端的时期。20世纪30年代，我们遭遇了过去600年来最严重的大萧条。而今天几乎是过去600年来最兴旺的时期……我们使用了“悲惨繁荣”这个词语来描绘它，因为它产生的后果是悲惨的。
\par

债券比股票理性得多。没有人会以为债券的价值会上升到月球去。
\par

\subsubsection*{未来的市场回报率}

（2000年11月）如果股市未来的回报率能够达到15\%，那肯定是因为一种强大的“伦勃朗效应”。这不是什么好事。看看日本以前的情况就知道了。当年日本股市的市盈率高达50～60倍。这导致了长达十年的经济衰退。不过我认为日本的情况比较特殊。我的猜测是，美国不会出现极端的“伦勃朗化”现象，今后的回报率将会是6\%。如果我错了，那可能是因为一个糟糕的原因。如果未来股票交易更像伦勃朗作品的买卖，那么股市将会上涨，但缺乏坚实的支撑。如果出现这种情况，那就很难预测它会发展到什么程度，股价将会上涨多高，以及会持续到什么时候。
\par

（2001年4月）总的来说，我认为美国投资者应该调低他们的预期。人们不是有点蠢，而是非常蠢。没有人愿意说出这一点。
\par

\begin{WisdomBox}
\WisdomSpeaker{查理·芒格}
\par

人们不是有点蠢，而是非常蠢。没有人愿意说出这一点。
\par
\end{WisdomBox}

（2001年5月）我们并没有处在一个猪猡也能赚钱的美好时代。投资游戏的竞争变得越来越激烈。
\par

我不认为我们将会看见大多数股票像1973～1974年，甚至是1982年那样，价格跌到令人流口水的低位。我觉得沃伦和我非常有可能再也看不到这些机会了，但那并不是坏事。我们只要继续努力就好啦。
\par

（巴菲特：不过也不是完全没有可能。你无法预测市场将会怎样。在日本，十年期债券的利率只有0.625\%。以前谁能想得到呢？）
\par

就算日本发生了那样的事情，美国的情况也不会那么糟糕。在接下来这段日子里，那些普通的高薪投资顾问将不会有非常好的表现。
\par

垃圾债的短暂崩溃（2002年底）是很奇怪的事情，当时它们的收益率高达35\%～40\%。垃圾债市场绝对是一片混乱。如果股市也遭遇这种情况——那么不难想象，很快就会大崩盘。
\par

\subsubsection*{当前的市场环境}

（2004年5月）说到大环境，我认为它会令那些想要很容易、很有把握地赚到钱的人感到非常痛苦。普通股的估值可能相当合理，但它们不算特别便宜。
\par

我们的现金就足以说明问题了。如果我们有很多好主意，我们就不会拥有这么多现金。
\par

伯克希尔和西科有大量我们不知道用来干什么的现金。如果你算上债券的话，伯克希尔有700亿美元，西科的现金也是多得不得了。我们的现金是有史以来最多的。要是在以前，我们只要沉住气就可以了，我们能够把现金派上用场。
\par

\subsection*{华尔街出卖了道德}

\subsubsection*{华尔街的道德缺失}

华尔街的平均道德水准永远至多只能说是中等……这并不意味着华尔街没有一些正直的聪明人，其实是有的，像今天在座的各位，但我认识的每个人都不得不和他自己的公司对抗（为了做正确的事情）。
\par

过去这些年来，投资银行业的道德风气逐渐败坏。许多年前，我们花600万美元收购多元零售公司（Diversified Retailing），银行对我们进行了严格而聪明的审查。那些银行家关心他们的客户，想要保护他们的客户。
\par

现在的风气是，只要能赚钱，什么都可以卖。你能把东西卖掉，你就道德高尚。但这个标准可不够。
\par

\begin{WisdomBox}
\WisdomSpeaker{查理·芒格}
\par

（华尔街）现在的风气是，只要能赚钱，什么都可以卖。你能把东西卖掉，你就道德高尚。但这个标准可不够。
\par
\end{WisdomBox}

\subsubsection*{所罗门个案研究}

现在的情况真令人吃惊。所罗门的自律和理性程度至少不比其他投资银行差，但到头来，所罗门居然求爷爷告奶奶地想要得到罗伯特·麦克斯韦尔这种人的投资业务，这人的绰号可是“跳来跳去的捷克人”啊。他既然有这个花名，投资银行就不应该去抢着做他的生意。
\par

［巴菲特：就在人们发现他跳海的那天(由于丑行败露，他跳海自杀了)，我们(所罗门)把钱汇给他，以交换他正要汇给我们的钱，但他没有付款。所以，我们去英国找他的几个儿子收账，那真是一团糟。我们收回了应得的钱。］
\par

如果有个投资银行家多开几张交易单给麦克斯韦尔，那么他的收入将会显著增长。你得防止人们为了挣钱将不三不四的人带入。
\par

\subsubsection*{诺曼底美国的个案研究}

（在谈起所罗门和一家他们没有点名的不良企业有交易时，沃伦和查理自顾自地笑了起来。然而，在2002年的年会上，查理点明那家企业就是诺曼底美国公司。所罗门发现这家公司公开招股的材料全是伪造的，所以在尚未发生交易的时候及时地取消了它的首次公开募股。根据美国证券交易委员会网站上的一份报告：“1995年8月15日，诺曼底的股票被推荐在纳斯达克全国市场系统上交易。一天之后，诺曼底取消了新股发行，并宣布所有的交易无效。”）
\par

这件事是一个很好的（又充满讽刺的）经历。沃伦、路易斯·辛普森和我都是（所罗门的）董事，我们还是最大的股东，我们说：“别跟这个家伙做生意。”但他们不听我们的话，还说承销委员会已经批准了那次新股发行。
\par

［沃伦：“他身上挂着一个霓虹灯牌写着‘骗子’。他确实坐牢了。碰巧的是，他宣称拥有大量伯克希尔的股票，并从中获利不菲，但我检查了股东名录，并没有看到他的名字。他的名字也许很普通，跟他重名的人很多，但如果他真有那么多股份，我想我应该能够找出来（所以他很可能在持有伯克希尔股份这件事上说了谎）。”］
\par

\subsection*{芒格对金融机构和衍生品的警告}

\subsubsection*{金融机构的风险}

金融机构的本质决定了它很容易阴沟里翻船。你有可能放贷太多，进行了愚蠢的收购，或者过度借贷——不仅仅是衍生品（才会让你倒闭）。
\par

也许没有人像我们这样，但我们对金融风险特别敏感。当金融机构想努力表现时，我们就会感到紧张。
\par

我们尤其害怕那些大量举债的金融机构。如果他们开始说起风险管理有多么好，我们就会很紧张。
\par

我们比其他人更加警惕。如果感到不安，我们很快就会把大量的钱抽出来。这是我们的经营方式——你必须适应它。
\par

\subsubsection*{衍生品}

衍生品系统简直是神经病，它完全不负责任。人们以为的固定资产并不是真正的固定资产。它太复杂了，我在这里没办法说清楚——但你无法相信涉及的金额达到几万亿美元那么多。你无法相信它有多么复杂。你无法相信衍生品的会计工作有多难。你无法相信人们对企业的价值和清算能力的一厢情愿的想法会产生多大的激励作用。
\par

人们不去考虑后果的后果：人们开始只想对冲利率变化带来的风险，这是非常困难和复杂的事情。然后呢，对冲操作导致了（公布的盈利）波动。所以他们用新的衍生品来抹平。这个时候你就不得不说谎了。这就变成了“疯帽匠的茶话会”。那些庞大的、复杂的公司会出现这种情况。
\par

必须有人站出来说：“我们不做这种事情——因为这太难了。”
\par

我认为，要测试任何（有显著衍生品风险的）大型机构的管理人员的智力和道德水准，最好的方法是问他们：“你真的能理解你的衍生品账目吗？”如果有人说“是”，那这个人要么疯掉了，要么是在说谎。
\par

\begin{WisdomBox}
\WisdomSpeaker{查理·芒格}
\par

要测试任何（有显著衍生品风险的）大型机构的管理人员的智力和道德水准，最好的方法是问他们：“你真的能理解你的衍生品账目吗？”如果有人说“是”，那这个人要么疯掉了，要么是在说谎。
\par
\end{WisdomBox}

说起能源衍生品（的情况），你会更容易理解（这种危险）——它们全部完蛋了。当（那些公司）想要动用他们的账本上写着的财产时，钱并不在那里。而说到金融资产，这样的结局还没有出现，会计规则也还没变，所以这种结局会在将来出现。
\par

衍生品有很多条款说，如果一方的信用被调低，那么它必须增加抵押金。那就像股票账户的贷款保证金一样——你将会（由于增加更多的保证金而）破产。为了保护自己，他们引入了不稳定的因素。似乎没有人意识到他们创造的系统是一个大灾难。那是一个疯狂的系统。
\par

在工程领域，人们拥有很大的安全边际。但在金融界，人们根本不考虑安全。他们只是追求扩张，扩张，扩张。做假账为此推波助澜。对此我比沃伦更加悲观。
\par

\subsubsection*{衍生品的会计法}

我非常讨厌将GAAP（公认会计准则）应用于衍生品和掉期协议。摩根大通为了虚报的前端收入将自己彻底出卖给了这种会计准则。我觉得这种做法很无耻。
\par

这种做法很疯狂，会计人员出卖了自己的职业操守。每个人都屈服了，采用了宽松的（会计）标准，创造出和各种理论模型相关的怪异衍生品。结果就是，产生了大量由会计人员捏造出来的利润，而不是真实的收益。当你想动用那笔钱的时候，它就消失不见了。它从来就不曾存在过。
\par

它（衍生品的会计法）就是这么恶心。它像阴沟那么臭，如果我猜得没错，那些公司将会付出沉重的代价。你所能做的就是想好当账本上那些衍生品破灭时该怎么办。
\par

说美国衍生品的会计法像阴沟是对阴沟的侮辱。
\par

\begin{WisdomBox}
\WisdomSpeaker{查理·芒格}
\par

说美国衍生品的会计法像阴沟是对阴沟的侮辱。
\par
\end{WisdomBox}

\subsubsection*{衍生品破灭的可能性}

我们试图卖掉通用再保险公司的衍生品业务，但没有成功，所以我们开始清算它。我们不得不大幅削减价格。我非常有把握地预测（美国的）主要银行无法以账本上的价格清算它们持有的衍生品。结局什么时候出现，会有多么糟糕，我现在说不准。但我担心后果会很可怕。我认为将会出现比能源领域更大的问题，你看看能源领域被破坏成什么样就知道了。
\par

在未来五到十年里，如果美国没有遇到（和衍生品相关的）大问题，我将会非常吃惊。
\par

我认为我们是美国唯一甩掉衍生品账本的大型公司。
\par

对于已经很有钱的公司而言，比如说伯克希尔，去从事这种业务是很神经的想法。对大银行来讲，衍生品是一桩疯狂的生意。
\par

\begin{WisdomBox}
\WisdomSpeaker{查理·芒格}
\par

对大银行来讲，衍生品是一桩疯狂的生意。
\par
\end{WisdomBox}

\subsection*{会计业的堕落遗患无穷}

\subsubsection*{大型会计师事务所的道德沦丧}

在我年轻的时候，大型的会计师事务所是非常正派的地方，没有人赚脏钱。但在过去的25年来，它们逐渐做出一些糟糕的行为。你一旦开始做坏事，那么第二次做就容易多了，到最后，你就会成为道德败坏的人。这个国家的大型会计师事务所居然为了钱而明目张胆地做假账，帮助客户逃税……
\par

太多的律师和会计师事务所参与了不可告人的事情。例如，它们收取应变费和保密费，提供逃税服务，这绝对是令人厌恶的事情。
\par

\begin{WisdomBox}
\WisdomSpeaker{查理·芒格}
\par

太多的律师和会计师事务所参与了不可告人的事情。
\par
\end{WisdomBox}

为了让你更好地理解那些顶尖的伪造避税用的财务报表的审计师事务所坏到什么程度，我想告诉你一件事，有个人曾经跟我说，他们比其他事务所做得更好，因为他们只为最大的20家客户提供这项服务，所以没有人会发现。
\par

\subsubsection*{激进会计法}

我们被（在美国企业界泛滥的）激进会计法吓坏了，所以尽量采用保守的会计法。保守的会计法有助于我们的生意决策和保护伯克希尔。我们怎么可能去做那些可能有悖于我们的道德准则的事情呢？
\par

创造性会计（伪造账目）绝对是对文明的诅咒。人们可以说复式簿记法是历史的一大进步。可是利用会计来造假和骗人是可耻的行为。在民主国家里，往往需要丑闻才能撬动改革。安然是企业文化长期出毛病的最典型例子。
\par

\subsubsection*{EBITDA的误用}

我认为，你每次看到“EBITDA”（即未计利息、税项、折旧及摊销前的利润）这个词，你都应该用“狗屁利润”来代替它。
\par

\subsubsection*{安达信}

我认为（所发生的一切对安达信那些无辜职员来讲）是非常不公平的，但没有破产的资本主义，就像没有地狱的宗教。当事情糟糕到这种地步，又缺乏内部控制体系——安达信就没有（内控体系）——也许应该有个事务所要倒掉。
\par

\begin{WisdomBox}
\WisdomSpeaker{查理·芒格}
\par

但没有破产的资本主义，就像没有地狱的宗教。
\par
\end{WisdomBox}

我们不会做任何有可能导致伯克希尔垮台的事情。安达信特别脆弱，因为它是合伙制企业。合伙制企业在采取举动和选择客户等方面必须额外谨慎。
\par

\subsubsection*{美国养老基金的会计丑闻}

当前这种（假设养老基金能够获得高得脱离现实的回报率）做法是愚蠢而且不正当的。但如果你跟管理人员提出来，他们的眼睛在敌意来临之前就已经瞪上你了。
\par

养老基金的会计由于利用了脱离现实的假设而陷入了丑闻。以最近的事情推断未来是人类的本性，但要是公司管理者们也这么想，那就太糟糕了。
\par

IBM刚刚将它对养老基金回报率的预期调高到10\%。绝大多数公司（在2000年）的回报率是9\%。我们认为6\%更为现实一些。该公司的领导也许相信这个数字是可以达到的——他们都是诚实的人——但从潜意识层面来说，他们相信它，是因为他们想要相信它。它能让盈利显得很好，这样他们就能推动股价上涨了。
\par

下面这句谚语很好地总结了那些会计人员什么话也不说的原因：“谁给我面包吃，我就给谁唱歌。”我认为你正从美国的会计系统得到一些非常愚蠢的数据。我觉得他们可能不是故意要撒谎，但也有可能是。
\par

\subsubsection*{糟糕的会计法导致道德沦丧}

有着宽松会计标准的人会引来其他人做出极其可怕的举动。那是造孽，绝对是造孽。如果你携带几大桶钞票经过贫民窟，又不好好看管它们，那么你就犯了很严重的罪行，因为你会惹来许多糟糕的行为，而糟糕的行为是会传染的。同样地，命令会计师草率地做账是一种真正的罪孽，也是一种愚蠢的经商方法。
\par

\begin{WisdomBox}
\WisdomSpeaker{查理·芒格}
\par

有着宽松会计标准的人会引来其他人做出极其可怕的举动。那是造孽，绝对是造孽。
\par
\end{WisdomBox}

\subsection*{芒格对股票期权的批评}

\subsubsection*{作为报酬的股票期权}

如果你瞧瞧股票期权的后果，你将会看到许多可怕的行为。给一个亲手创办公司、已经六十几岁的CEO大量的期权，试图以此刺激他对该公司的忠诚，这是一种精神错乱的想法。玛约医疗中心（Mayo Clinic）的医生或者克拉法斯律所（Cravath）的律师在六十几岁时还会为了期权而更加努力地工作吗？
\par

我们的股东也知道，我们的体系不同于绝大多数大型的公司。我们认为它不那么变化无常。股票期权体系可能会使某些什么事情都没做的人得到大量的回报，而应该得到很多报酬的人却什么都没有。除非我们收购的公司原来就实行了（股票期权制度），否则我们不会使用它。
\par

\subsubsection*{期权估值和布莱克—舒尔斯模型}

布莱克—舒尔斯模型对短期的期权而言是有效的，但如果期权是长期的，而且你认为你（对期权涉及的资产）有所了解，那么使用布莱克—舒尔斯模型是一种很神经的做法。
\par

布莱克—舒尔斯模型是一套“一无所知”的系统。如果你对价值毫不了解，只了解价格，那么利用布莱克—舒尔斯模型，你能够很准确地猜出为期90天的期权到底值多少钱。但如果期权比90天更长久，那么使用这个模型就有点神经啦。
\par

举例说吧，我们在好市多发行了执行价分别为30美元和60美元的股票期权，而布莱克—舒尔斯模型认为60美元的期权价值更高。这简直是疯了。
\par

\subsubsection*{股票期权的会计法}

那种认为期权没有成本的理论导致大量的公司滥用期权，这种情况对美国来说是很糟糕的，因为人们会认为公司的薪酬是不公平的。
\par

我已经厌倦了这个话题。这个话题我已经讲了很久啦。用错误的会计法来经营一个文明世界实在是很糟糕的。这就好比修建桥梁的时候，把工程设计搞错了。当那些颇有声望的人也说股票期权不应该算为成本支出的话，这就太过分了。
\par

股票期权不但会增加企业的成本，还会稀释股权。任何反对这种看法的人都是神经病。
\par

\begin{WisdomBox}
\WisdomSpeaker{查理·芒格}
\par

股票期权不但会增加企业的成本，还会稀释股权。任何反对这种看法的人都是神经病。
\par
\end{WisdomBox}

这些人（投票支持不将股票期权算为成本的88名参议员）是愚蠢而可耻的。他们知道那是错误的，偏偏还那么做。
\par

这件事从头到尾都让人恶心。
\par

我宁愿在妓院弹钢琴谋生，也不愿像约翰·多尔（注：著名硅谷风险投资家，但因强烈反对将股票期权算作成本而广受诟病）推荐的那样，靠股票期权来赚钱。
\par

\begin{WisdomBox}
\WisdomSpeaker{查理·芒格}
\par

我宁愿在妓院弹钢琴谋生，也不愿像约翰·多尔推荐的那样，靠股票期权来赚钱。
\par
\end{WisdomBox}

\subsection*{对基金管理企业的批评}

\subsubsection*{有缺陷的激励机制}

现在一般基金的管理系统要求人们假装去做他们做不了的事情，假装喜欢他们其实不喜欢的事情。这种生活方式很可怕，但它能得到很多薪水。
\par

\begin{WisdomBox}
\WisdomSpeaker{查理·芒格}
\par

现在一般基金的管理系统要求人们假装去做他们做不了的事情，假装喜欢他们其实不喜欢的事情。这种生活方式很可怕，但它能得到很多薪水。
\par
\end{WisdomBox}

\subsubsection*{没有增加价值}

开放式基金每年收取2\%的费用，然后那些经纪商让人们在不同基金之间转来转去，多付3\%～4\%的费用。可怜的普通基民们把钱交给专业人士，却得到糟糕的结果。我觉得这很恶心。本来人们买你的产品，你应该让他们的资产增值。但在这个国家，基金管理人员却只顾着自己赚钱。
\par

这个行业很滑稽，因为从净收益来看，整个投资管理行业加起来并没有对全部客户的资产提供附加值。这就是这个行业的运作方式。
\par

\begin{WisdomBox}
\WisdomSpeaker{查理·芒格}
\par

这个行业很滑稽，因为从净收益来看，整个投资管理行业加起来并没有对全部客户的资产提供附加值。
\par
\end{WisdomBox}

\subsubsection*{股票经纪人vs指数基金}

这次年会坐满了许多聪明的、值得尊敬的股票经纪人，他们为客户提供了很好的服务，在这里很难指责他们。但总的来说，股票经纪人的表现太差了，甚至还比不上指数基金。
\par

我认为对于一般基金来说，投资指数是比现在通行的无借贷股票投资更为聪明的投资策略——尤其是在现行做法每年的管理费用超过1\%的情况下。指数基金的表现不可能永远都这么好，如果几乎所有人都转而投资它的话。但它的良好表现可以持续很长一般时间。
\par

\subsubsection*{共同基金丑闻}

选择投资管理人员这件事最近变得更难了，因为人们发现，有不少共同基金的管理人员接受贿赂，出卖了他们自己股东的利益。
\par

这就像一个人走过来说：“我们为什么不把你母亲杀死，然后瓜分她的保险金呢？”很多人说：“好啊，我正需要那笔保险金呢。”
\par

他们有许多人觉得自己被冤枉了。
\par

如果共同基金的监事是独立的，那么我就是俄罗斯波尔修芭蕾舞团的首席男明星。
\par
